\chapter*{\uppercase{ABSTRACT}}
\addcontentsline{toc}{section}{\bfseries \uppercase{Abstract}}
\begin{spacing}{1.5}
\begin{sloppypar}
%Stroke is the leading cause of mortality and disability worldwide, making the rapid and precise identification of stroke types—normal, ischemic, or hemorrhagic—essential for effective clinical decision-making. Although CT scans are the most commonly used method for initial evaluations, automated classification of stroke types from these images faces difficulties due to limited datasets, class imbalance, and subtle differences in the appearance of the lesion. This study introduces a deep learning framework that employs generative AI for balancing the dataset, alongside multi-model training utilizing Vision Transformer (ViT), ConvNeXt, and CoAtNet architectures, coupled with explainable AI for improved visual understanding. The data set undergoes preprocessing and augmentation using a GAN-based method to tackle the issue of class imbalance. The interpretability of the model is enhanced by Grad-CAM++ and attention rollout visualizations, which emphasize key areas related to stroke lesions, thus promoting clinical trust and increasing model transparency. The proposed system delivers an accurate, interpretable, and comprehensive AI solution for reliable stroke detection and classification based on CT scans.


Traditional autoscaling mechanisms based on static CPU and memory thresholds often fail to handle sudden traffic spikes in microservices architectures, particularly in university course registration systems. This project proposes an intelligent autoscaling system using queuing theory for a distributed course registration platform. The system models each microservice as an M/M/c queue and uses real-time metrics like request arrival rate and service latency to compute performance indicators. 

A Multi-Feature Decision System (MFDS) determines scaling actions by comparing predicted response time against SLA thresholds, while a Minimum Scaling Overhead (MSO) strategy computes optimal replica counts in a single step. The implementation uses containerized microservices with Docker Swarm orchestration, Prometheus monitoring, and dynamic scaling integration. Experimental evaluation demonstrates faster response times, reduced scaling latency, and improved resource utilization compared to traditional threshold-based autoscaling methods.
\end{sloppypar}
\end{spacing}